
\keyheading{第1题　轴对称及对称轴}
\begingroup\color{answerred}三个图形都是轴对称图形；对称轴依次为1条、2条、4条。\par\noindent\begin{center}{%
\begin{tikzpicture}[x=1bp,y=-1bp,color=black,draw=black,text=black]
\begin{scope}\clip (65.00,133.00) rectangle (551.00,231.00);
\hdhead{76}{图形的运动——信息窗1}
\hdsection{1}{145}{下面的图形是轴对称图形吗？如果是，分别画出它们所有的}
\q{140}{188}{690}{13.5}{\color{bookred}对称轴。}
\hdsection{2}{361}{在（　　）里写出下面各图形对称轴的条数。}
\hdsection{3}{698}{画出下面每个图形的另一半，使它们成为轴对称图形。}
\begin{scope}[scale=.6666667]
\node[font=\normalsize,inner sep=0pt] at (121,241) {①};
\node[font=\normalsize,inner sep=0pt] at (375,241) {②};
\node[font=\normalsize,inner sep=0pt] at (625,240) {③};
\draw[hdoutline] (143.5,277)--(158.5,234)--(217.5,234)--(232.5,277)--cycle;
\draw[hdoutline] (405.5,239)--(440.5,260)--(475.5,239)--(475.5,309)--(440.5,288)--(405.5,309)--cycle;
\draw[hdoutline] (673.5,234)--(705.5,234)--(705.5,258)--(729.5,258)--(729.5,290)--(705.5,290)--(705.5,314)--(673.5,314)--(673.5,290)--(649.5,290)--(649.5,258)--(673.5,258)--cycle;
\node[font=\normalsize,inner sep=0pt] at (123,415) {①};
\node[font=\normalsize,inner sep=0pt] at (378,415) {②};
\node[font=\normalsize,inner sep=0pt] at (627,415) {③};
\node[font=\normalsize,inner sep=0pt] at (123,553) {④};
\node[font=\normalsize,inner sep=0pt] at (378,553) {⑤};
\node[font=\normalsize,inner sep=0pt] at (627,553) {⑥};
\draw[hdoutline] (147,432)--(234,432)--(234,480)--(147,480)--cycle;
\draw[hdoutline] (174,450)--(191,432)--(209,450)--cycle;
\draw[hdoutline] (175,450)--(208,450)--(208,480)--(175,480)--cycle;
\draw[hdoutline] (186,460)--(195,460)--(195,480)--(186,480)--cycle;
\draw[hdoutline] (158,419)--(174,419)--(174,432)--(158,432)--cycle;
\draw[hdoutline] (156,419)--(166,404)--(176,419)--cycle;
\draw[hdoutline] (163,440)--(168,440)--(168,448)--(163,448)--cycle;
\draw[hdoutline] (163,452)--(168,452)--(168,460)--(163,460)--cycle;
\draw[hdoutline] (163,464)--(168,464)--(168,472)--(163,472)--cycle;
\draw[hdoutline] (207,419)--(223,419)--(223,432)--(207,432)--cycle;
\draw[hdoutline] (205,419)--(215,404)--(225,419)--cycle;
\draw[hdoutline] (212,440)--(217,440)--(217,448)--(212,448)--cycle;
\draw[hdoutline] (212,452)--(217,452)--(217,460)--(212,460)--cycle;
\draw[hdoutline] (212,464)--(217,464)--(217,472)--(212,472)--cycle;
\draw[hdoutline] (424,408)--(464,408)--(485,444)--(464,480)--(424,480)--(403,444)--cycle;
\draw[hdoutline] (693,442) circle (37);
\draw[hdoutline] (693,442) circle (31);
\draw[hdoutline] (671,415)--(720,464)--(715,469)--(666,420)--cycle;
\draw[hdoutline] (153,545)--(168,545)--(168,560)--(153,560)--cycle;
\draw[hdoutline] (213,545)--(228,545)--(228,560)--(213,560)--cycle;
\draw[hdoutline] (168,560)--(183,560)--(183,575)--(168,575)--cycle;
\draw[hdoutline] (198,560)--(213,560)--(213,575)--(198,575)--cycle;
\draw[hdoutline] (183,575)--(198,575)--(198,590)--(183,590)--cycle;
\draw[hdoutline] (168,590)--(183,590)--(183,605)--(168,605)--cycle;
\draw[hdoutline] (198,590)--(213,590)--(213,605)--(198,605)--cycle;
\draw[hdoutline] (153,605)--(168,605)--(168,620)--(153,620)--cycle;
\draw[hdoutline] (213,605)--(228,605)--(228,620)--(213,620)--cycle;
\draw[hdoutline] (444,545)--(447.827,572.761)--(460.971,565.029)--(453.239,578.173)--(481,582)--(453.239,585.827)--(460.971,598.971)--(447.827,591.239)--(444,619)--(440.173,591.239)--(427.029,598.971)--(434.761,585.827)--(407,582)--(434.761,578.173)--(427.029,565.029)--(440.173,572.761)--cycle;
\draw[draw=black,line width=.45pt] (444,582)--(444,545);
\draw[draw=black,line width=.45pt] (444,582)--(460.971,565.029);
\draw[draw=black,line width=.45pt] (444,582)--(481,582);
\draw[draw=black,line width=.45pt] (444,582)--(460.971,598.971);
\draw[draw=black,line width=.45pt] (444,582)--(444,619);
\draw[draw=black,line width=.45pt] (444,582)--(427.029,598.971);
\draw[draw=black,line width=.45pt] (444,582)--(407,582);
\draw[draw=black,line width=.45pt] (444,582)--(427.029,565.029);
\draw[hdoutline] (694,582) circle (38);
\draw[hdoutline] (694,544)--(716.336,612.743)--(657.86,570.257)--(730.14,570.257)--(671.664,612.743)--cycle;
\end{scope}
\q{141}{502}{140}{12.5}{（　　）条}
\q{141}{641}{140}{12.5}{（　　）条}
\q{397}{502}{140}{12.5}{（　　）条}
\q{397}{641}{140}{12.5}{（　　）条}
\q{648}{502}{140}{12.5}{（　　）条}
\q{648}{641}{140}{12.5}{（　　）条}
\begin{scope}[scale=.6666667]
\draw[hdgrid] (155,742)--(155,1159);
\draw[hdgrid] (180,742)--(180,1159);
\draw[hdgrid] (205,742)--(205,1159);
\draw[hdgrid] (229,742)--(229,1159);
\draw[hdgrid] (253,742)--(253,1159);
\draw[hdgrid] (277,742)--(277,1159);
\draw[hdgrid] (301,742)--(301,1159);
\draw[hdgrid] (325,742)--(325,1159);
\draw[hdgrid] (349,742)--(349,1159);
\draw[hdgrid] (373,742)--(373,1159);
\draw[hdgrid] (396,742)--(396,1159);
\draw[hdgrid] (420,742)--(420,1159);
\draw[hdgrid] (443,742)--(443,1159);
\draw[hdgrid] (467,742)--(467,1159);
\draw[hdgrid] (491,742)--(491,1159);
\draw[hdgrid] (514,742)--(514,1159);
\draw[hdgrid] (538,742)--(538,1159);
\draw[hdgrid] (562,742)--(562,1159);
\draw[hdgrid] (585,742)--(585,1159);
\draw[hdgrid] (608,742)--(608,1159);
\draw[hdgrid] (632,742)--(632,1159);
\draw[hdgrid] (655,742)--(655,1159);
\draw[hdgrid] (679,742)--(679,1159);
\draw[hdgrid] (703,742)--(703,1159);
\draw[hdgrid] (726,742)--(726,1159);
\draw[hdgrid] (750,742)--(750,1159);
\draw[hdgrid] (774,742)--(774,1159);
\draw[hdgrid] (155,742)--(774,742);
\draw[hdgrid] (155,767)--(774,767);
\draw[hdgrid] (155,791)--(774,791);
\draw[hdgrid] (155,815)--(774,815);
\draw[hdgrid] (155,840)--(774,840);
\draw[hdgrid] (155,865)--(774,865);
\draw[hdgrid] (155,889)--(774,889);
\draw[hdgrid] (155,914)--(774,914);
\draw[hdgrid] (155,938)--(774,938);
\draw[hdgrid] (155,963)--(774,963);
\draw[hdgrid] (155,988)--(774,988);
\draw[hdgrid] (155,1013)--(774,1013);
\draw[hdgrid] (155,1037)--(774,1037);
\draw[hdgrid] (155,1062)--(774,1062);
\draw[hdgrid] (155,1087)--(774,1087);
\draw[hdgrid] (155,1111)--(774,1111);
\draw[hdgrid] (155,1136)--(774,1136);
\draw[hdgrid] (155,1159)--(774,1159);
\draw[draw=black,dashed,line width=.55pt] (253,743)--(253,913);
\draw[draw=black,dashed,line width=.55pt] (288,918)--(288,1157);
\draw[draw=black,dashed,line width=.55pt] (585,743)--(585,963);
\draw[draw=black,dashed,line width=.55pt] (608,966)--(608,1158);
\draw[hdoutline] (205,815)--(253,767);
\draw[hdoutline] (585,767)--(632,816)--(632,891)--(585,939);
\draw[hdoutline] (608,990)--(632,1038)--(679,1062)--(632,1087)--(608,1133);
\end{scope}
\q{160}{1165}{598}{8}{原扫描左侧局部轮廓缺失，仅重绘可辨认的线段；左下图不能唯一恢复。}

\begin{scope}[color=answerred,draw=answerred,text=answerred]\ans{104.67}{196.00}{104.67}{10}{是，1条}
\ans{271.33}{216.00}{104.00}{10}{是，2条}
\ans{434.00}{218.67}{109.33}{10}{是，4条}
\draw[answerred,dashed,line width=.8pt] (125.33,149.33) -- (125.33,191.33);
\draw[answerred,dashed,line width=.8pt] (293.67,152.67) -- (293.67,211.33);
\draw[answerred,dashed,line width=.8pt] (264.67,182.67) -- (322.67,182.67);
\draw[answerred,dashed,line width=.8pt] (459.67,150.00) -- (459.67,215.33);
\draw[answerred,dashed,line width=.8pt] (427.33,182.67) -- (491.33,182.67);
\draw[answerred,dashed,line width=.8pt] (434.67,158.00) -- (484.67,208.00);
\draw[answerred,dashed,line width=.8pt] (434.67,208.00) -- (484.67,158.00);
\ans{116.67}{334.67}{20.67}{12}{$1$}
\ans{286.67}{334.67}{20.67}{12}{$6$}
\ans{453.33}{334.67}{20.67}{12}{$2$}
\ans{116.67}{426.67}{20.67}{12}{$4$}
\ans{286.67}{426.67}{20.67}{12}{$4$}
\ans{453.33}{426.67}{20.67}{12}{$5$}
\draw[answerred,line width=1pt] (168.67,511.33) -- (201.33,543.33);
\draw[answerred,line width=1pt] (390.00,511.33) -- (358.67,544.00) -- (358.67,593.33) -- (390.00,626.00);
\draw[answerred,line width=1pt] (405.33,660.00) -- (389.33,692.00) -- (357.33,708.00) -- (389.33,724.67) -- (405.33,755.33);
\ans{206.00}{560.00}{138.67}{8.5}{可辨斜边已补。}
\ans{107.33}{650.67}{234.00}{9.5}{左下原图轮廓不清，\\无法唯一确定待补半图。}
\end{scope}
\end{scope}\pgfresetboundingbox\path[use as bounding box] (65.00,133.00) rectangle (551.00,231.00);
\end{tikzpicture}}
\end{center}\par\endgroup
\keyheading{第2题　对称轴的条数}
\begingroup\color{answerred}\keygrid{3}{
$1$ & $6$ & $2$\\
$4$ & $4$ & $5$\\
}\noindent 单位：条。\par\endgroup
\keyheading{第3题　补成轴对称图形}
\begingroup\color{answerred}\textbf{原扫描限制：}左侧部分轮廓不清，不能唯一补全；仅补已确认线段。右侧两图答案如下。\par\smallskip\noindent\begin{center}\begin{adjustbox}{max width=.68\linewidth}\begin{tikzpicture}[x=.48cm,y=.48cm,color=answerred]\draw[step=1,answerred!25,thin] (-3,-1) grid (3,8);\draw[dashed] (0,-1)--(0,8);\draw[black,thick] (0,7)--(2,5)--(2,2)--(0,0);\draw[answerred,very thick] (0,7)--(-2,5)--(-2,2)--(0,0);\node[below] at (0,-1){右上图};\begin{scope}[xshift=5cm]\draw[step=1,answerred!25,thin] (-3,-1) grid (3,7);\draw[dashed] (0,-1)--(0,7);\draw[black,thick] (0,6)--(1,4)--(3,3)--(1,2)--(0,0);\draw[answerred,very thick] (0,6)--(-1,4)--(-3,3)--(-1,2)--(0,0);\node[below] at (0,-1){右下图};\end{scope}\end{tikzpicture}\end{adjustbox}\end{center}\par\endgroup

\begingroup\color{answerred}\textbf{多解核对：}对折法和对应点法应给出相同的对称轴及补图结果；不能用“答案不唯一”代替原扫描缺失信息。\par\endgroup
